\begin{textsample}{POS Dim 1 – pe – Score 35.00 – pe/pe\_ef\_71.txt}  \label{ex:f1_pos_250}
6.4 HISTÓRIA O conhecimento histórico \textbf{emerge} como \textbf{disciplina} com pretensões científicas no \textbf{século} XIX.
Associada, em grande medida, ao aparecimento dos Estados \textbf{Modernos} e arrogando -se a responsabilidade de contar a história desses Estados e, por \textbf{conseguinte}, da formação da nação e do povo, necessários à constituição de cada um deles, ficava, portanto, reservada à História a função de estabelecer os marcos, os feitos, os fatos e o panteão de heróis que teriam sido responsáveis pela história da nação e da construção de uma dada identidade nacional.
Na \textbf{análise} de Elza Nadai : > O \textbf{século} XIX acrescentou, paralelamente aos grandes movimentos que ocorreram visando construir os Estados Nacionais sob hegemonia burguesa, a necessidade de retornar -se ao passado, com o objetivo de identificar a “ base comum ” formadora da nacionalidade.
Daí os conceitos tão caros às histórias nacionais : Nação, Pátria, Nacionalidade, Cidadania ( NADAI, 1986, p. 106 ).
Nesse sentido, a História tinha por finalidade forjar, mediante o estabelecimento de uma continuidade entre passado-presente-futuro, os traços da nacionalidade nos indivíduos de um dado território, proporcionando -lhes a construção de uma memória e de uma tradição.
Ela buscava assumir, assim, não só uma função científica na ordem do saber, mas também um sentido magisterial na formação dos sujeitos necessários aos \textbf{ideais} nacionalistas de composição dos Estados \textbf{Modernos}.
É nesse contexto que a História aparecerá, pela primeira vez, para o pensamento ocidental, como uma \textbf{disciplina} que precisa ser ensinada.
A História era pensada como mestra da vida.
Conhecer e aprender com a história tinha um claro sentido exemplar : fazer com que as gerações futuras não repetissem os \textbf{erros} do passado e aprendessem com os exemplos dos grandes homens em seus feitos heroicos e exemplares.
No entanto, essa modalidade de escrita e ensino de História trazia embutida em seus \textbf{fundamentos} um claro viés elitista de legitimação dos estratos superiores da sociedade, ligados ao establishment dos Estados \textbf{Modernos}, além de servir de narrativa de justificação de uma dada realidade política e social e sua reprodução.
Grosso modo, essa concepção de História se traduzia, reproduzia e circulava mediante uma prática de ensino positivista, tradicional, que ecoava uma história das e para as elites políticas, alçadas à condição de únicos sujeitos produtores da História.
Não tardou e no próprio \textbf{século} XIX ocorreram as primeiras críticas a essa modalidade de escrita da história e às finalidades ideológicas de sua reprodução.
Ainda naquele período, o filósofo alemão Karl Marx colocou -se como um \textbf{crítico} voraz daquela concepção e a acusou de estar a serviço de uma ideologia aristocrática, burguesa e excludente.
O marxismo propôs uma escrita da história politicamente engajada e comprometida com a conscientização da classe trabalhadora.
Nesse sentido, o ensino da História deveria servir à conscientização das massas para a revolução.
A história é \textbf{investida} não só de um sentido teleológico, mas também é pensada como orientadora da ação política para a transformação da realidade social e revolução da própria história, mediante o estabelecimento de uma nova ordem, de um novo tempo.
O ensino de história teria, em seu \textbf{horizonte}, a formação e a conscientização da classe trabalhadora, forjando um sujeito \textbf{crítico} e revolucionário que toma consciência de que é sujeito de sua própria história.
O marxismo retira das elites o privilégio de se \textbf{dizerem} os únicos fazedores da história.
Marx defende que “ os homens fazem a história, mesmo que não o saibam ”.
Seria função dos historiadores e/ou dos professores de história, como intelectuais de vanguarda, promover, mediante o processo de \textbf{ensino-aprendizagem}, essa tomada de consciência do sujeito histórico, fazendo -o \textbf{crítico} e autônomo para que pudesse agir e intervir politicamente sobre sua realidade e modificá -la ( REIS, 1999, 2003 ).
Mas, foi da Escola dos Annales, no início do \textbf{século} XX, que vieram as maiores alterações \textbf{tanto} na produção do conhecimento histórico e da sua escrita quanto nas implicações para o seu ensino.
É com os Annales que a História vai ser pensada, talvez pela primeira vez, em relação às demais \textbf{ciências} sociais.
Problematizada, de acordo com Marc Bloch, como “ a \textbf{ciência} dos homens no tempo ”, a História vai ser pensada na sua \textbf{interface} com a Antropologia, a Sociologia, a Psicologia, a Economia, a \textbf{Ciência} Política, etc.
Esse diálogo possibilita uma ampliação dos \textbf{métodos}, das técnicas, sobretudo dos \textbf{problemas}, dos objetos, das \textbf{abordagens} e, em especial, das fontes necessárias para sua escrita.
Isso amplia de forma considerável a dimensão e o quantitativo dos sujeitos que ganham dignidade historiográfica, ou seja, que passam a ser vistos não só como sujeitos históricos, mas também como produtores de narrativas sobre o passado.
Há, dessa \textbf{maneira}, uma quebra em relação aos modos de se produzir, conhecer e ensinar História estabelecidos, até então, sobretudo em relação à tradição dita positivista e sua concepção linear, teleológica e elitista de História ( REIS, 1999, 2003 ).
Essa tendência à pluralização e à diversificação dos estudos históricos se acentuou ainda mais a partir da segunda metade do \textbf{século} XX, com o desenvolvimento de \textbf{inúmeras} outras modalidades de produção e escrita da História.
Essas foram decorrentes da crise de paradigmas que atravessou as Ciências Humanas a partir daquele período e que questionava, radicalmente, as grandes narrativas que buscavam dar sentido à realidade e à própria história da \textbf{humanidade}.
A História não é mais vista como mestra da vida.
Ela se \textbf{investe} de outras funções cada vez mais específicas e contextualizadas a dados períodos e espaços.
Muito embora preserve uma função magisterial bem clara : “ fazer \textbf{lembrar} o que uma dada sociedade quer esquecer ”.
Dessa forma, a produção do conhecimento histórico vai ser diretamente relacionada aos \textbf{problemas} do presente.
O historiador é aquele que conhece o passado a partir dos \textbf{problemas} colocados no e pelo presente.
Ele passa a se pensar \textbf{tanto} como sujeito quanto como objeto do próprio conhecimento que produz em uma relação permanente de tensão entre presente e passado na qual este é constantemente reescrito à medida que, como defende Paul Veyne ( 1992 ), o repertório de questões muda ou se amplia.
É nesse sentido que o ensino de História passa a ser pensado como promotor das diferenças, problematizador de identidades cristalizadas e direcionado à promoção da cidadania, da \textbf{pluralidade} e da democratização dos modos de vida.
Diante dessa perspectiva, cabe destacar o pensamento de Schmidt e Cainelli : > É importante destacar que a História tem uma função didática de formar uma consciência histórica cada vez mais \textbf{complexa}, com a perspectiva de fornecer elementos para a orientação, interpretação do passado, para dentro, [ des]construindo identidades, e para fora, fornecendo sentidos para a ação na vida prática ( SCHIMIDT e CAINELLI, 2009, p. 37 ).
É, portanto, diante dessas perspectivas e contextos de produção do conhecimento histórico que os currículos de História, \textbf{tanto} em nível nacional quanto em níveis estaduais e municipais, começaram a ser pensados no Brasil, sobretudo a partir da década de 1980.
Ora privilegiando uma dada \textbf{abordagem}, ora contemplando uma diversidade delas.
Prevista desde a LDBEN ( 1996 ), a construção de uma Base Nacional Comum Curricular foi objeto de discussão de educadores e especialistas ao longo das décadas seguintes, o que produziu uma série de documentos orientadores que vão dos \textbf{Parâmetros} Curriculares Nacionais, passando pelas Diretrizes Curriculares Nacionais até chegar às discussões que \textbf{fomentaram} a \textbf{atual} BNCC.
É nesse bojo de discussões que \textbf{emerge} a BNCC de História, trazendo em si todas estas tensões e disputas \textbf{inerentes} ao próprio campo do saber histórico, mas apontando para a necessidade de um curriculum mínimo que garanta o direito à aprendizagem em História de dados conteúdos, conceitos, práticas e o desenvolvimento de \textbf{competências} e habilidades que expressem a constituição de um dado sujeito, pensado como cidadão formador de uma sociedade justa, democrática e plural ( BNCC, 2017 ).
Esse desafio proposto pela BNCC coloca para os Estados alguns \textbf{problemas}, em especial para a construção do curriculum de História, uma vez que esses terão de compatibilizar as especificidades de suas realidades sociais, culturais e históricas com um conjunto de \textbf{competências} e habilidades gerais já estabelecidas pelo documento nacional a partir de uma \textbf{abordagem}, sobretudo nos anos \textbf{finais} do ensino fundamental, cronológica.
Apesar de estabelecer que esse é apenas um arranjo possível e que os Estados e municípios podem fazer ou construir outros, a BNCC dificulta a possibilidade de pensá -los a partir de eixos temáticos, como é feito nos “ \textbf{Parâmetros} Curriculares de História do Estado de Pernambuco ” ( PERNAMBUCO, 2013 ).
Dessa \textbf{maneira}, o currículo que se seguirá será pensado levando em consideração essas variadas \textbf{abordagens}, sem desconsiderar os entendimentos diversos, procurando deixar margem de flexibilidade para que os municípios e demais sistemas de ensino possam \textbf{imprimir} as características que \textbf{julgarem} necessárias quando este documento chegar ao chão da sala de \textbf{aula}.

% matched lemmas: abordagem, análise, atual, aula, ciência, competência, complexo, conseguinte, crítico, disciplina, dizer, emergir, ensino-aprendizagem, erro, final, fomentar, fundamento, horizonte, humanidade, ideal, imprimir, inerente, interface, investir, inúmero, julgar, lembrar, maneira, moderno, método, parâmetro, pluralidade, problema, século, tanto
\end{textsample}
